\newif\ifijgisclass
\IfFileExists{tGIS2e.cls}{\ijgisclasstrue}{\ijgisclassfalse}

\ifijgisclass
  \documentclass{tGIS2e}
\else
  \documentclass[11pt,a4paper]{article}
\fi

\ifijgisclass
  \usepackage{graphicx}
  \usepackage{booktabs}
  \usepackage{amsmath,amssymb}
  \usepackage{multirow}
  \usepackage{url}
  \usepackage{algorithm}
  \usepackage{algorithmic}
  \usepackage{subcaption}
  \citestyle{tGIS}
\else
  \usepackage[margin=1in]{geometry}
  \usepackage[T1]{fontenc}
  \usepackage[utf8]{inputenc}
  \usepackage{lmodern}
  \usepackage{graphicx}
  \usepackage{booktabs}
  \usepackage{amsmath,amssymb}
  \usepackage{multirow}
  \usepackage{url}
  \usepackage[numbers,sort&compress]{natbib}
  \usepackage{algorithm}
  \usepackage{algorithmic}
  \usepackage{subcaption}
  \usepackage[hidelinks]{hyperref}
  \newcommand{\name}[1]{#1}
  \newcommand{\affil}[1]{#1}
  \newcommand{\correspondence}[1]{\par\noindent\textbf{Correspondence:} #1}
  \newenvironment{keywords}{\par\noindent\textbf{Keywords: }}{\par}
  \newcommand{\jvol}[1]{}
  \newcommand{\jnum}[1]{}
  \newcommand{\jyear}[1]{}
  \newcommand{\jmonth}[1]{}
  \newcommand{\DOI}[1]{}
  \newcommand{\received}[1]{}
  \newcommand{\revised}[1]{}
  \newcommand{\accepted}[1]{}
  \newcommand{\pubyear}[1]{}
  \newcommand{\pagerange}[1]{}
\fi

\usepackage{placeins}
\graphicspath{{../../experiments/outputs/figures/}}

\title{A dynamic rock--TBM spatial relation graph for component-indexed geological anomaly readouts during TBM excavation}

\author{
\name{First Author\textsuperscript{a}\thanks{Email: first.author@example.com} and
Second Author\textsuperscript{b}} \\
\affil{\textsuperscript{a}Department of XXX, University A, City, Country; \\
\textsuperscript{b}Institute of YYY, University B, City, Country}
}

\jvol{00}
\jnum{00}
\jyear{2026}
\jmonth{Month}
\DOI{00.0000/0000000000000000}
\received{00 Month 2026}
\revised{00 Month 2026}
\accepted{00 Month 2026}
\pubyear{2026}
\pagerange{1--20}

\begin{document}

\maketitle

\begin{abstract}
The coupling between underground geological space and mobile engineering equipment is a typical spatiotemporal process in mechanized tunnelling. In TBM excavation, construction anomalies are difficult to interpret from monitoring responses alone because hazardous or weak ground affects the machine through specific spatial relations with the cutterhead, shields, and other components. Modern TBMs record dense streams of thrust, torque, penetration, pressure, and other responses, while tunnel seismic prediction (TSP) describes where anomalous ground may occur. However, these data do not directly answer the key spatial question: which geological units are geometrically related to which TBM components as excavation advances. This study proposes a dynamic rock–TBM spatial relation graph to provide this missing relation layer. The graph represents TSP-derived rock voxels and component-labelled TBM surface samples as two types of spatial entities, and connects them only when active-zone, distance, facing-direction, and excavation-state constraints indicate a plausible rock–machine relation. Rebuilding the graph at each excavation step converts static geological anomalies into component-specific geological anomaly trajectories along the tunnel alignment. These trajectories provide a spatial context for interpreting monitoring residuals and can be traced back to the rock voxels, surface samples, and relation attributes that produce them. Two tunnel cases demonstrate how the graph supports excavation anomaly interpretation beyond monitoring-only or geological-field-only summaries. In the SJLS case, the cutterhead-indexed trajectory remains strongly associated with shield-pressure residuals after detrending, whereas pooling all TBM surface components removes this relation pattern. Its stepwise trajectory is also consistent with shield-pressure residual changes, indicating that the graph captures not only where anomalous ground exists, but how its component-specific relation to the machine evolves during excavation. The BSLL case further verifies that the same graph construction supports component-level querying and edge-level traceability under 3D voxel support, although its compact intervals provide weaker residual evidence. The proposed graph is therefore a queryable spatial-relation layer for TBM excavation analysis and future safe-construction workflows, rather than a contact-force measurement, calibrated risk model, or universal response predictor.

\end{abstract}

\begin{keywords}
spatial representation; rock--machine relation graph; anomaly readout field;
candidate relation; persistence residual; TBM excavation
\end{keywords}

\ifijgisclass\else
\correspondence{First Author. Email: first.author@example.com}
\fi

\section{Introduction}
Tunnel boring machines (TBMs) are widely used in long and deep tunnel construction,
where excavation responses emerge from the coupled effects of geological conditions,
machine geometry, and operational process. Under fractured, water-rich, high-stress, or
locally deformable ground conditions, the interaction state between surrounding rock
and the TBM may change rapidly and manifest through abnormal thrust, torque,
penetration, advance rate, or shield-related responses. These monitoring variables are
informative external expressions of rock--TBM interaction, but the geological body and
the cutterhead or shield form a spatial relation that evolves as excavation advances.
Monitoring curves alone provide limited evidence about where the interaction occurs,
which machine components are implicated, and how local interaction relations evolve
along chainage.

Most existing methods process TSP-derived geological indicators and TBM parameters as
tabular variables or time-series variables. Such representations can support trend
prediction, but they cannot answer which rock spatial unit forms a candidate interaction
relation with which machine component at a given excavation step
\citep{yagiz2008utilizing,hassanpour2010tbm,armaghani2018prediction,hochreiter1997long}.
Sequence models preserve temporal dependence but treat excavation as an ordered signal
rather than a spatially structured interaction process, so rock location, machine
component structure, and local rock--machine candidate relations are not retained within
a unified representation.

Rather than pursuing a more complex prediction model, this study constructs a
chainage-referenced rock--TBM interaction graph sequence that organises rock voxels,
TBM surface samples, and operational response sequences into a computable
spatio-temporal graph. TSP-derived geological attributes are formalised as rock voxel
nodes, while the cutterhead and shield are discretised as parameterised TBM surface
nodes. Candidate rock--machine edges are screened through spatial proximity,
signed-normal consistency, active-zone filtering, and excavation-state constraints.
Instead of interpreting edge weights as contact-force measurements or relying on
sparse jamming labels, the framework uses continuously recorded monitoring responses to
supervise the relative importance of geometry-screened candidate relations, and
aggregates the learned relevance at edge, surface-node, component, and chainage
levels to generate response-consistent diagnostic views.

The main contributions are threefold. First, the study proposes a chainage-referenced
rock--TBM spatial representation that formalises heterogeneous spatial entities---TSP
voxels, TBM surface components, and chainage-indexed monitoring responses---into a
unified entity model preserving geometric and topological properties. Second, it
proposes a geometry-constrained rock--machine interaction graph sequence in which
candidate edges are restricted by active-zone membership, distance, surface-normal
compatibility, and excavation state, ensuring that learned relations remain
geometrically plausible rather than unconstrained. Third, it uses response-supervised
prediction to examine the response consistency of the learned graph relations and
outputs edge-, component-, and chainage-level relevance diagnostics, with clear
interpretation limits, as spatially organised diagnostics rather than contact-force or
risk estimates.


\section{Related work}
\label{sec:related_work}

\subsection{TBM response modelling and spatial gaps}
Traditional TBM performance studies estimate penetration, advance rate,
thrust, torque, or related operational responses from rock mass properties and
machine parameters. Empirical, statistical, and machine-learning formulations
are used for planning and performance assessment in hard-rock tunnelling
\citep{yagiz2008utilizing,hassanpour2010tbm,armaghani2018prediction}. With the
accumulation of monitoring data, sequence models such as LSTM and GRU also
learn response inertia and lagged dependence from recent operational histories
\citep{hochreiter1997long,cho2014learning}. These studies provide important
response-prediction baselines, but their spatial unit is usually a tunnel
section, a sample, or a feature vector rather than an explicit relation between
a geological object and a TBM component.

This limitation is central to the present paper. A monitoring curve can
describe how torque, thrust, penetration, advance rate, or shield pressure
changes with time or chainage, but it does not state which geological voxels
are plausible interaction partners for the cutterhead, front shield, middle
shield, or tail shield. The present study therefore does not try to replace
monitoring-based predictors. It asks whether a spatial relation representation
provides component-indexed field information about the rock--machine
configuration that monitoring curves alone do not encode.

\subsection{Geological prediction as spatial support}
Advance geological prediction methods, including TSP and related geophysical
or drilling-based approaches, are used to identify faults, fractured zones,
water-rich zones, weak surrounding rock, and other geological anomalies before
or during tunnel excavation \citep{li2019tunnel}. For TBM analysis, such data
are valuable because they describe spatial heterogeneity that is not contained
in the monitoring record itself. TSP results can be represented as velocity
profiles, sections, or voxelised attribute fields along the tunnel direction;
voxelisation is useful here because it preserves spatial support, local
attribute values, and chainage-referenced geological context.

Geological spatial representation alone is nevertheless incomplete for the
diagnostic question addressed in this paper. A voxel field can show where
anomalous ground occurs, but it does not specify whether that anomaly is
adjacent to the cutterhead or shield, whether the adjacency is
normal-compatible, or whether the voxel remains active at the current
excavation step. This is a spatial-support problem as much as an engineering
problem: TSP voxels, TBM surface samples, and monitoring records have
different spatial or observational supports and cannot be interpreted as a
single tabular unit without losing relation information
\citep{goodchild1992geographical,longley2005geographic,openshaw1984modifiable}.

\subsection{Entity--relation spatial data models}
GIScience is concerned not only with mapping locations but also with
representing geographical objects, spatial relations, temporal change, and
processes \citep{goodchild1992geographical,worboys2004gis,galton2001spatial}.
Spatial relation theory provides concepts for neighbourhood, adjacency,
containment, intersection, and other topological or metric relations
\citep{egenhofer1991point,wu2008contiguity}. TBM excavation fits this
perspective because rock units, machine surfaces, and monitoring responses are
observed in a moving excavation coordinate, and the relevant relations change
as the machine advances.

Recent geospatial knowledge-graph and graph data-model studies show how
heterogeneous spatial entities are integrated for retrieval, reasoning, and
query \citep{zhu2026knowwheregraph,xu2025spatiotemporal,li2025question,rahimi2021topology}.
Scenario-oriented spatiotemporal knowledge-graph models likewise emphasise
linking heterogeneous data through meaningful objects and relations rather
than storing them only as parallel indicator tables
\citep{wang2025urbanphysical}. The present paper shares this entity--relation
modelling perspective, but it operates at a finer voxel and surface-node
support and targets a moving underground engineering process rather than an
indoor social network, an urban knowledge service, or a regional knowledge
base.

This literature provides the conceptual foundation for the proposed graph:
rock voxels are geological objects, TBM surface samples are machine-side
spatial objects, and rock--machine edges encode screened candidate interaction
relations. Monitoring responses are kept as aligned records rather than graph
nodes. The representational gap is therefore not the absence of a predictor,
but the lack of a spatial data model that connects geological voxels, machine
components, candidate interaction relations, and response residuals in one
component-indexed field representation.

\subsection{Spatial interaction graphs}
Graphs are a natural formalism for representing irregular spatial objects and
their relations. Recent geospatial studies have used graph-based
representations for spatial classification, spatial multivariate
distributions, urban network morphology, and spatially constrained community
detection
\citep{zhu2020spatial,desabbata2023graph,liu2025urban,chen2026usgt}. Spatial
interaction models further emphasise that relations are shaped by distance
decay, compatibility, direction, accessibility, and spatial support
\citep{wilson1971spatial,fotheringham1989spatial}. For the rock--TBM problem,
this means that candidate relations should not be fully connected or purely
learned without physical constraints. They should depend on active-zone
membership, distance, surface-normal compatibility, and excavation state.

Graph neural networks and heterogeneous graph models remain relevant as future
extensions for learning on structured data
\citep{schlichtkrull2018modeling,zhang2019heterogeneous,wu2020comprehensive,zhou2020graph}.
Spatiotemporal graph models have also been successful in traffic, flow
prediction, and physical simulation
\citep{li2018diffusion,wu2019graph,zhang2020deep,geng2019spatiotemporal,sanchez2020learning,pfaff2021learning}.
However, these studies mainly motivate the representational capacity of graph
structures and the need for spatially constrained relations. They do not imply
that the present tunnel cases should be framed primarily as a black-box
graph-learning benchmark.

In the tunnelling domain, digital twin frameworks for TBM performance
prediction, visualisation, monitoring, and deformation prediction have been
proposed \citep{latif2023digital,apoji2023shaping,hu2025dataknowledge}. The
present graph-based representation is complementary: rather than replacing
prediction-oriented digital twin modules, it specifies an underground spatial
relation layer between geological voxels and TBM surface components. This layer
defines active candidate relations, component-indexed aggregation, and
edge-level traceability, which are typically implicit in monitoring-centred
digital twin workflows. The evaluation is therefore kept descriptive:
Spearman association, detrending, null comparisons, sensitivity analysis, and
edge-level traceability are used to support a spatial field-representation
claim, not a universal prediction or risk-calibration claim.


\section{Methodology}
\label{sec:methods}
This section formalises rock--TBM interaction as an excavation-step-indexed
spatial entity-relation graph. The method starts from a TSP-derived rock voxel
field, an advancing parameterized TBM surface, and aligned monitoring records
under a common excavation-step index. It then constructs a sparse rock--TBM
interaction graph at each excavation step using active-zone, distance,
normal-compatibility, and excavation-state constraints. Finally, weighted
rock--machine edges are aggregated by TBM component to compute a component
$\times$ excavation-step anomalous interaction field. This field supports
excavation anomaly indication, component-indexed readout, and
jamming-risk-oriented interpretation while retaining traceable edge
contributions.
The graph is sparse because rock--machine edges are not fully connected; they
are retained only when the four spatial and excavation-state constraints are
jointly satisfied.
Figure~\ref{fig:pipeline} summarises the method structure.

The three subsections correspond to the input representation and two core
algorithmic operations. Section~\ref{sec:definition} defines the two graph
entity types and the aligned monitoring records. Section~\ref{sec:construction}
defines geometry-constrained sparse graph construction, and
Section~\ref{sec:interpretable_computation} defines component-level anomalous
interaction field computation from screened rock--machine edges.

The formulation estimates a queryable field rather than a direct
mechanical response. Given geological voxels $D_{\mathrm{geo}}$, a
parameterized TBM surface $M_{\mathrm{TBM}}$, and monitoring records aligned by
excavation step, the graph-construction operator
\[
\mathcal{G}: (D_{\mathrm{geo}},M_{\mathrm{TBM}},t)\mapsto G_t
\]
produces sparse rock--machine candidate relations, and the field-computation
operator
\[
\mathcal{F}: (G_t,\phi(g_i),\rho_j)\mapsto
\{I_c(t),\Delta I_c(t)\}_{c\in C,t=1,\ldots,T}
\]
aggregates those relations into a component-indexed anomalous interaction
field. The aligned monitoring records are used to read this field against
excavation responses; they are not inputs to $\mathcal{G}$ and do not define
graph edges.

\begin{figure}[htbp]
\centering
\includegraphics[width=0.95\linewidth]{fig1_method_framework.pdf}
\caption{Method overview of the geometry-constrained rock--TBM spatial
interaction graph computation. (a) Sparse graph construction from active rock
voxels and component-labelled TBM surface nodes using active-zone, distance,
normal-compatibility, and excavation-state constraints. (b) TSP-derived rock
voxel field with geological attribute support. (c) Parameterized TBM surface
components advanced along the excavation direction. (d) Component $\times$
excavation-step anomalous interaction field computed by aggregating weighted
rock--machine edges by TBM component; the field can subsequently be aligned
with monitoring responses and queried back to contributing edges.}
\label{fig:pipeline}
\end{figure}

\subsection{Rock, machine, and monitoring data representation}
\label{sec:definition}
The input data are organised into two graph entity types and aligned
monitoring records: rock voxels, TBM surface nodes, and excavation monitoring
responses aligned to excavation steps. Monitoring records are not graph nodes;
they are used after graph computation to evaluate and interpret the
component-level anomalous interaction field.

\paragraph{Rock voxel.} The geological input is a TSP-derived voxel field
$D_{\mathrm{geo}}=\{(c_i,g_i)\}$, where $c_i=(x_i,y_i,z_i)$ is the coordinate
of rock voxel $i$ and $g_i$ is its geological attribute vector (e.g.\
P-wave velocity $V_p$, S-wave velocity $V_s$, dynamic Young's modulus, velocity
ratio, Poisson's ratio, density). Rock voxels are the geological spatial
support of the graph.

\paragraph{TBM surface node.} The machine input is a parameterized TBM
surface. Each surface node is represented by $(p_j(0),n_j,\rho_j)$, where
$p_j(0)$ is the initial surface-node coordinate, $n_j$ is the local outward
surface normal, and $\rho_j$ is the machine component label. The current
implementation partitions the machine surface into cutterhead, front shield,
middle shield, and tail shield components; $\rho_j$ therefore indexes one of
these four labels. TBM surface nodes are the machine-side spatial support of
the graph.

\paragraph{Aligned monitoring records.} At each excavation step $t$ the monitoring
vector is
\[
u_t = [\mathrm{AdvanceRate}_t,\; \mathrm{RPM}_t,\; \mathrm{Torque}_t,\;
\mathrm{Thrust}_t,\; \mathrm{Penetration}_t,\; \mathrm{ShieldPressure}_t],
\]
and the response variables used downstream are denoted $r_t^{(k)}$, where $k$
indexes a response channel such as torque, thrust, penetration, advance rate,
or shield pressure. Monitoring responses are aligned to the graph output
$I_c(t)$ and $\Delta I_c(t)$; they do not participate in graph construction.

\paragraph{Common chainage coordinate.} Because TSP coordinates and TBM
monitoring mileage are commonly recorded in different local frames, the two
graph entity types and the monitoring records are mapped to a common local
chainage coordinate. With
$\mathrm{ShieldMileage}(t)$ the shield mileage stamp at step $t$, a simple
alignment is
\[
x_{\mathrm{local}}(t)=
\mathrm{ShieldMileage}(t)-\mathrm{ShieldMileage}_{\mathrm{start}},
\]
after which TBM surface nodes are advanced along the excavation direction
$d$:
\[
p_j(t+1)=p_j(t)+\Delta x\,d .
\]
Rock voxels, TBM surface nodes, and monitoring records are therefore aligned
to the same excavation-step index $t$. This subsection only establishes that
common index; it does not treat monitoring records as graph entities.

Table~\ref{tab:spatial_entities_relations} summarises the resulting
entity--relation data model. The table makes explicit that the graph is a queryable
entity-relation representation rather than only a plotting device or a
neural-network input tensor. Figure~\ref{fig:pipeline}(b--d) illustrates how
the rock voxel field, TBM surface components, and component-indexed field
readouts are formalised under the common excavation-step reference.

\begin{table}[htbp]
\centering
\caption{Entities, relations, and readouts in the geometry-constrained rock--TBM interaction
graph computation.}
\label{tab:spatial_entities_relations}
\scriptsize
\setlength{\tabcolsep}{3pt}
\begin{tabular}{p{0.17\linewidth}p{0.43\linewidth}p{0.30\linewidth}}
\toprule
Entity or quantity & Definition & Role \\
\midrule
$V_t^r$ & Active TSP-derived rock voxels & Geological spatial support \\
$V_t^m$ & Parameterized TBM surface samples & TBM component support \\
$E_t^{rr}$ & Rock--rock voxel neighbourhoods & Geological contiguity context \\
$E_t^{mm}$ & TBM surface neighbourhoods & Machine-surface topology context \\
$E_t^{rm}$ & Screened rock--machine candidate relations & Candidate interaction support \\
$A_c(t)$ & Sum of geometry weights for component $c$ & Auxiliary exposure readout \\
$I_c(t)$ & Geometry-weighted anomaly intensity for component $c$ & Main field readout \\
$e_{t+h}^{(k)}$ & Persistence residual of response $k$ & Aligned monitoring readout \\
\bottomrule
\end{tabular}
\end{table}

\subsection{Geometry-constrained sparse graph construction}
\label{sec:construction}
Algorithm~\ref{alg:graph_construction} summarises the graph construction
operation. Its purpose is to produce a sparse and queryable relation support
$E_t^{rm}$ for each excavation step, rather than to infer measured contact.

\begin{algorithm}[htbp]
\caption{Geometry-constrained sparse graph construction}
\label{alg:graph_construction}
\begin{algorithmic}[1]
\REQUIRE Rock voxels $\{(c_i,g_i)\}$, TBM surface nodes $\{p_j(0),n_j,\rho_j\}$, excavation steps $t=1,\ldots,T$
\ENSURE Graph snapshots $G_t=(V_t^r\cup V_t^m,E_t^{rr}\cup E_t^{mm}\cup E_t^{rm})$
\FOR{each excavation step $t$}
\STATE Advance the TBM surface to $p_j(t)$ using the common excavation-step coordinate.
\STATE Select active rock voxels $V_t^r$ using the face and shield active zones.
\STATE Construct $V_t^m$, $E_t^{rr}$, and $E_t^{mm}$ from voxel and TBM-surface neighbourhoods.
\FOR{each active rock voxel $i$ and TBM surface node $j$}
\IF{$d_{ij}(t)\leq \tau_{\mathrm{edge}}$, $\kappa_{ij}(t)\geq \eta_{\min}$, and $s_i(t)=1$}
\STATE Add $(i,j)$ to $E_t^{rm}$ and assign $w_{ij}^{rm}(t)=\exp[-d_{ij}(t)/\tau_{\mathrm{edge}}]\kappa_{ij}(t)$.
\ENDIF
\ENDFOR
\ENDFOR
\end{algorithmic}
\end{algorithm}

At each excavation step a graph snapshot is defined as
\[
G_t=\left(V_t^r \cup V_t^m,\; E_t^{rr} \cup E_t^{mm} \cup E_t^{rm}\right),
\]
where $V_t^r$ and $V_t^m$ are the rock and TBM surface node sets defined in
Section~\ref{sec:definition}, $E_t^{rr}$ denotes rock--rock voxel adjacency,
$E_t^{mm}$ denotes TBM surface adjacency, and $E_t^{rm}$ denotes candidate
rock--machine relations. Rock node features are
\[
x_i^r(t)=[c_i,g_i,s_i(t)],
\]
where $s_i(t)$ records whether voxel $i$ remains active and eligible for
interaction at step $t$. TBM node features are
\[
x_j^m(t)=[p_j(t),n_j(t),\rho_j].
\]
The intra-rock edge set $E_t^{rr}$ is constructed from voxel neighbourhoods,
and $E_t^{mm}$ is constructed from local neighbourhoods on the parameterized
TBM surface. The central graph-design decision is the construction of
$E_t^{rm}$, the rock--machine candidate edges.

\paragraph{Edge-screening rules.} A candidate edge $(i,j)$ is included in
$E_t^{rm}$ only when four conditions are jointly satisfied:
\[
\begin{cases}
c_i\in\Omega_t,\\
d_{ij}(t)\leq \tau_{\mathrm{edge}},\\
\kappa_{ij}(t)\geq \eta_{\min},\\
s_i(t)=1,
\end{cases}
\]
namely active-zone membership, a distance threshold, surface-normal
compatibility, and an excavation-state constraint. The first rule restricts
candidate relations to the excavation-relevant rock volume $\Omega_t$; the
second uses
\[
d_{ij}(t)=\|c_i-p_j(t)\|
\]
to keep only spatially close rock--machine pairs; the third uses the signed
surface-normal compatibility
\[
\kappa_{ij}(t)=
\max\left(0,\frac{n_j(t)^\top(c_i-p_j(t))}{d_{ij}(t)+\epsilon}\right)
\]
to keep only geometrically plausible facing relations; and the fourth prevents
already excavated voxels from repeatedly contributing to later graph snapshots.

The active zone $\Omega_t$ separates face-contact plausibility from
shield-side proximity within one active rock volume:
\[
\Omega_t=\Omega_t^{\mathrm{face}}\cup\Omega_t^{\mathrm{shield}} ,
\]
with, using cutterhead chainage $x_f(t)$, tail-shield chainage
$x_{\mathrm{tail}}(t)$, radial distance $r_i=\sqrt{y_i^2+z_i^2}$, cutterhead
radius $R_c$, shield radius $R_s$, face look-ahead length $L_f$, and radial
active-zone buffer $\tau_{\mathrm{zone}}$,
\[
\Omega_t^{\mathrm{face}}=
\left\{i \mid 0 \leq x_i-x_f(t) \leq L_f,\;
r_i \leq R_c+\tau_{\mathrm{zone}}\right\},
\]
\[
\Omega_t^{\mathrm{shield}}=
\left\{i \mid x_{\mathrm{tail}}(t)\leq x_i\leq x_f(t),\;
R_s \leq r_i \leq R_s+\tau_{\mathrm{zone}}\right\}.
\]
Table~\ref{tab:active_zone_parameters} lists the parameter values used in all
formal case runs. The radii and shield length come from the parameterized TBM
geometry. The edge-screening parameters $\tau_{\mathrm{edge}}$ and
$\eta_{\min}$ are varied in Section~\ref{sec:sensitivity_settings}. The
active-zone parameters $L_f$ and $\tau_{\mathrm{zone}}$ define the main
spatial support and are fixed in the reported formal runs.

\begin{table}[htbp]
\centering
\caption{Active-zone and candidate-relation parameters used in all formal
case runs.}
\label{tab:active_zone_parameters}
\scriptsize
\setlength{\tabcolsep}{5pt}
\begin{tabular}{lll}
\toprule
Parameter & Value & Role \\
\midrule
$L_f$ & 5.0\,m & Face look-ahead length for $\Omega_t^{\mathrm{face}}$ \\
$\tau_{\mathrm{zone}}$ & 5.0\,m & Radial active-zone buffer around face and shield \\
$R_c$ & 4.0\,m & Parameterized cutterhead radius \\
$R_s$ & 3.95\,m & Parameterized shield radius \\
Shield length & 10.0\,m & Axial support for $\Omega_t^{\mathrm{shield}}$ \\
$\tau_{\mathrm{edge}}$ & 2.0\,m & Candidate-edge distance threshold in the main setting \\
$\eta_{\min}$ & 0.3 & Minimum surface-normal compatibility in the main setting \\
Surface sampling & 0.5\,m & TBM surface-node sampling resolution \\
\bottomrule
\end{tabular}
\end{table}

\begin{figure}[htbp]
\centering
\includegraphics[width=0.95\linewidth]{fig3_geometry_constrained_edges.pdf}
\caption{Geometry-constrained rock--machine candidate relation construction.
Candidate relations are controlled by active-zone membership, distance decay,
surface-normal compatibility, and excavation state. The construction defines
spatially plausible relations; it does not recover measured contact.}
\label{fig:geometry_edges}
\end{figure}

\paragraph{Edge weight.} For each candidate relation the geometry-compatibility
weight is
\[
w_{ij}^{rm}(t)=
\exp\left(-\frac{d_{ij}(t)}{\tau_{\mathrm{edge}}}\right)\kappa_{ij}(t).
\]
The weight increases when a rock voxel is closer to the TBM surface and more
compatible with the local surface normal. It is not a contact force or
pressure. With $V_t^r$, $V_t^m$, $E_t^{rr}$, $E_t^{mm}$,
$E_t^{rm}$, and $w_{ij}^{rm}(t)$ defined, the graph entities of
Section~\ref{sec:definition} have become a computable rock--TBM interaction
graph.

\subsection{Component-indexed anomalous interaction field computation}
\label{sec:interpretable_computation}
Once the graph is constructed, the central graph computation is to aggregate
screened rock--machine edges into a component-level anomalous interaction
field. Response alignment and edge-level query are downstream readouts of this
field rather than separate method branches.

\begin{algorithm}[htbp]
\caption{Component-level anomalous interaction field computation}
\label{alg:interaction_field}
\begin{algorithmic}[1]
\REQUIRE Graph snapshots $G_t$, anomaly scores $q_i=\phi(g_i)$, component labels $\rho_j$
\ENSURE Component $\times$ excavation-step field $\{I_c(t),\Delta I_c(t)\}$
\FOR{each excavation step $t$}
\FOR{each TBM component $c$}
\STATE Select rock--machine candidate edges $E_{c,t}^{rm}=\{(i,j)\in E_t^{rm}\mid \rho_j=c\}$.
\STATE Compute geometric exposure $A_c(t)=\sum_{(i,j)\in E_{c,t}^{rm}}w_{ij}^{rm}(t)$.
\STATE Compute anomaly intensity $I_c(t)=\sum_{(i,j)\in E_{c,t}^{rm}}w_{ij}^{rm}(t)q_i/(A_c(t)+\epsilon)$.
\STATE Compute step change $\Delta I_c(t)=I_c(t)-I_c(t-1)$ when $t>1$.
\ENDFOR
\ENDFOR
\STATE Store edge-level contributions $s_{ij}(t)=w_{ij}^{rm}(t)q_i$ for subsequent query.
\end{algorithmic}
\end{algorithm}

\paragraph{Component-level field.} The geological contribution of each rock
voxel is represented by a restrained TSP-derived anomaly score
\[
q_i=\phi(g_i).
\]
In the main formulation $q_i$ is a bounded low-velocity anomaly score derived
from the training active-zone distribution of P-wave velocity:
\[
q_i=
\operatorname{clip}\left(
\frac{Q_{95}^{\mathrm{train}}(Vp)-Vp_i}
{Q_{95}^{\mathrm{train}}(Vp)-Q_{5}^{\mathrm{train}}(Vp)+\epsilon},
0,1
\right).
\]
This fixed training-reference scaling avoids per-chainage min--max
normalisation and keeps field values comparable along chainage; lower
$V_p$ values correspond to higher geological anomaly scores under a common
reference distribution.

For each TBM component $c$, let $V_c^m$ denote the set of TBM surface nodes
whose component label is $c$. Two component-level quantities are defined at
each excavation step.
The pure geometric exposure is
\[
A_c(t)=
\sum_{j\in V_c^m}
\sum_{i:(i,j)\in E_t^{rm}}
w_{ij}^{rm}(t),
\]
which reports the geometric exposure of component $c$ to the screened
candidate rock--machine relation set. The geology-weighted spatial interaction
intensity is
\[
I_c(t)=
\frac{
\sum_{j\in V_c^m}
\sum_{i:(i,j)\in E_t^{rm}}
w_{ij}^{rm}(t)q_i
}{
\sum_{j\in V_c^m}
\sum_{i:(i,j)\in E_t^{rm}}
w_{ij}^{rm}(t)+\epsilon
},
\]
which reports the weighted geological anomaly intensity in the screened
candidate relations incident to component $c$.
$I_c(t)$ is defined for cutterhead, front shield, middle shield, and tail
shield at each chainage step. Together with its step-level change
$\Delta I_c(t)$, it forms the component $\times$ excavation-step anomalous
interaction field. $A_c(t)$ is retained as an auxiliary exposure readout, not
as part of the anomalous interaction field itself. Here, the field denotes the
component $\times$ excavation-step matrix of $I_c(t)$ and $\Delta I_c(t)$
values, whereas a field readout denotes a downstream use of that matrix through
response alignment or edge-contribution query.

\paragraph{Step-level tracking.} Beyond the instantaneous intensity
$I_c(t)$, the field also supports a step-level change view
\[
\Delta I_c(t) = I_c(t) - I_c(t-1),
\]
which reports how the anomalous interaction intensity around component $c$
changes between consecutive excavation steps. Tracking $I_c(t)$ and
$\Delta I_c(t)$ over the test chainage interval gives a component-indexed
anomaly trajectory that indicates where and when the rock--machine
interaction intensity rises or falls along the alignment.

The computed field is subsequently aligned with monitoring residuals for
response readout and can be queried back to edge contributions. For response
alignment, short-horizon TBM responses are evaluated through persistence
residuals rather than raw responses:
\[
e_{t+h}^{(k)} = r_{t+h}^{(k)}-r_t^{(k)},
\]
where $h$ is the diagnostic horizon and $k$ indexes a response variable. The
diagnostic question is whether component-indexed field readouts co-vary with
the response change left after persistence has been removed:
\[
I_c(t)\rightarrow e_{t+h}^{(k)} .
\]
The reported association statistic is the Spearman correlation between
$I_c(t)$ and $e_{t+h}^{(k)}$ over the evaluated test chainage interval. This
readout links the graph-computed field to TBM response change; it is
interpreted as field--residual co-variation within the evaluated chainage
interval, not as causal attribution. The choice of diagnostic pairs, null
variants, detrending, and robustness checks is part of the case evaluation
design and is specified in Section~\ref{sec:diagnostic_readout_variants}.
For edge-level traceability, candidate edges for any fixed diagnostic sample
are ranked by their contribution
\[
w_{ij}^{rm}(t)\,q_i,
\]
and each contributing edge retains the rock voxel identity, the TBM surface
node identity, the component label $\rho_j$, the distance $d_{ij}(t)$, the
normal compatibility $\kappa_{ij}(t)$, the geometry weight $w_{ij}^{rm}(t)$,
the anomaly score $q_i$, and the weighted contribution. The graph therefore
not only produces a component-indexed field but also traces that field back to
specific rock--machine spatial relations.

\FloatBarrier


\section{Case data and evaluation design}
\subsection{Case study data}
The experiments use two real tunnel cases. The first is the BSLL tunnel case
starting from DyK1017+205. It provides the main graph-sequence construction
example, the one-step response prediction run, and the three-step
advance-prediction sensitivity run. The second is the SJLS tunnel case starting
from Dyk1252+411. It uses an external TSP-derived velocity field and real
mileage-aligned TBM monitoring labels, and is used to test whether the same
spatial representation remains meaningful in a second tunnel setting.

The two cases are not treated as sequential data releases or replacement
datasets. They are two case-study settings with different evidential roles.
BSLL is used mainly to illustrate multi-step prediction and spatial
interpretation under a compact test interval. SJLS is used mainly for
prediction consistency and geometry-ablation evidence under an external TSP
case.

\begin{figure}[htbp]
\centering
\includegraphics[width=0.9\linewidth]{fig2_graph_construction.pdf}
\caption{Schematic of geometry-constrained graph construction at a
representative excavation step (illustrative, not drawn from the actual TSP
velocity field). The framework links TSP-derived rock voxels, parameterized TBM
surface nodes, and geometry-screened rock--machine candidate edges before
learning response-consistent relevance.}
\label{fig:tsp_profile}
\end{figure}

\subsection{Monitoring and geological variables}
For each excavation step, the monitoring vector contains advance rate, cutterhead
RPM, torque, thrust, penetration, and shield pressure. The prediction target is
the future multivariate response vector consisting of advance rate, torque,
thrust, penetration, and shield pressure. These variables describe the observed
machine response to the surrounding rock and are used both as historical inputs
and as response-supervision targets.

The geological input is a TSP-derived voxel field. For BSLL, the project TSP
voxel field contains P-wave velocity, S-wave velocity, density, dynamic Young's
modulus, velocity ratio, and Poisson's ratio around the tunnel alignment. For
SJLS, dense external P-wave and S-wave velocity profiles at Dyk1252+411 are
preprocessed into the same project TSP schema and paired with mileage-aligned
monitoring labels. This preprocessing allows both tunnel cases to use the same
graph-construction and graph-sequence learning pipeline.

\subsection{Mileage-ordered partitioning}
All experiments use mileage-ordered partitioning to preserve forward prediction.
For a historical window of $K=5$, each sample uses the preceding monitoring and
graph snapshots as input and predicts a future response at horizon $h$. The
formal case runs are summarised in Table~\ref{tab:data_summary}. The BSLL h=1
run evaluates one-step prediction, the BSLL h=3 run evaluates multi-step
advance prediction, and the SJLS run evaluates the external TSP case.

\begin{table}[htbp]
\centering
\caption{Formal case runs used in the experiments. Effective samples are
sequence windows after applying the history window and prediction horizon.}
\label{tab:data_summary}
\begin{tabular}{lcccc}
\toprule
Case & Tunnel & Start chainage & Horizon & Train/Test samples \\
\midrule
BSLL h=1 & BSLL & DyK1017+205 & 1 & 30 / 8 \\
BSLL h=3 & BSLL & DyK1017+205 & 3 & 29 / 7 \\
SJLS h=1 & SJLS & Dyk1252+411 & 1 & 76 / 17 \\
\bottomrule
\end{tabular}
\end{table}


\section{Results}
\subsection{Scope and interpretive boundaries}
Before presenting the experimental results, three scope limitations
should be noted to avoid over-interpretation.  First, the learned
edge relevance and surface relevance $C_j$ reflect
response-consistent interaction patterns under the forecasting task;
they are \emph{not} direct measurements of physical contact force,
contact pressure, or jamming probability.  Second, the case study
uses a single tunnel section with 49 excavation steps (44 effective
sequence samples, 12 in the stratified test set); all quantitative
indicators should be read as illustrative rather than statistically
conclusive.  Third, the spatial consistency indicators (Moran's $I$,
component CV, attention--geology correlation) quantify the internal
and external plausibility of the learned representation, not its
engineering correctness.  These boundaries are revisited in the
Discussion.

\subsection{Experimental setup}
\label{sec:experimental_setup}
All structural and training hyperparameters follow the Implementation
details in Section~\ref{sec:methods} (summarised here for readability):
historical window $K=5$, forecasting horizon $h=1$; candidate edge
distance threshold $\tau_{\mathrm{edge}}=2.0$\,m; signed-normal
compatibility threshold $\eta_{\min}=0.3$; active-zone radial extent
$\tau_{\mathrm{zone}}=5.0$\,m; TBM surface resolution 0.5\,m (1352 nodes);
GNN encoder with 2 layers, hidden dimension 32, dropout 0.3; GRU hidden
dimension 64, 1 layer; AdamW optimiser, learning rate $10^{-3}$, weight
decay $10^{-2}$, batch size 32, early stopping patience 15.  To control
for geological distribution shift between train and test, samples are
partitioned using stratified random sampling based on TSP shear wave
velocity ($V_s$) quartiles, yielding four geological strata each split
70/15/15 into train/val/test.  All models are evaluated on the
held-out stratified test set using MAE, RMSE, $R^{2}$, and Pearson
correlation.

\subsection{Prediction performance}
Table~\ref{tab:prediction} reports the global prediction performance
(Figure~\ref{fig:prediction}).  The proposed framework achieves MAE of
0.417 and Pearson $r$ of 0.567, outperforming all baselines on MAE and
achieving the only positive $R^{2}$ (0.133).  The stratified split
controls for geological heterogeneity between train and test, allowing
the geometric prior and learned rock--machine interactions to
generalise; under the previous mileage split, all models exhibited
negative $R^{2}$ due to a severe distribution shift in $V_s$ between
the training and test intervals (shift $= 2.31$ standard deviations,
Kolmogorov--Smirnov $p < 0.001$).  The proposed framework's primary
contribution lies in providing spatially explicit
interaction structure that monitoring-only models cannot access---a
trade-off examined in Section~\ref{sec:tradeoff}.

\begin{table}[htbp]
\centering
\caption{Global prediction performance on the stratified test set
($n=12$ samples drawn from four geological strata defined by TSP shear
wave velocity quartiles, ensuring the test distribution matches the
training distribution).  The proposed framework achieves the lowest MAE
and the only positive $R^{2}$ among all models, indicating that the
geometric prior and learned rock--machine interactions generalise when
geological heterogeneity between train and test is controlled.}
\label{tab:prediction}
\begin{tabular}{lcccc}
\toprule
Method & MAE & RMSE & $R^{2}$ & Pearson $r$ \\
\midrule
Persistence & 0.465 & 0.590 & $-0.008$ & 0.600 \\
XGBoost + TSP & 0.455 & 0.606 & $-0.114$ & 0.496 \\
LSTM & 0.488 & 0.602 & $-0.109$ & 0.457 \\
TSP-LSTM & 0.487 & 0.619 & $-0.174$ & 0.449 \\
Proposed framework & \textbf{0.417} & \textbf{0.540} & $\mathbf{0.133}$ & 0.567 \\
\bottomrule
\end{tabular}
\end{table}

\begin{figure}[htbp]
\centering
\includegraphics[width=0.9\linewidth]{fig3_prediction_comparison.pdf}
\caption{Predicted versus observed standardised monitoring responses
for five target variables on the stratified test set.  Each panel
corresponds to one response variable: (a)~advance rate, (b)~torque,
(c)~thrust, (d)~penetration, (e)~shield pressure.  Different curves
represent the proposed framework and baseline models; observed values
are shown as solid lines, predictions as dashed/dotted lines.  MAE
values in the legend are variable-specific values for the displayed
panel and should not be confused with the global multi-variable metrics
in Table~\ref{tab:prediction}.  Under the stratified split, the proposed
framework achieves positive $R^{2}$ and outperforms all baselines on
MAE, confirming that the geometric prior contributes predictive signal
when geological heterogeneity between train and test is controlled.}
\label{fig:prediction}
\end{figure}

\subsection{Structural ablation}
Table~\ref{tab:ablation} and Figure~\ref{fig:ablation} present the structural ablation results.  Four
ablation variants isolate the contribution of each design component.

\begin{table}[htbp]
\centering
\caption{Structural ablation results on the stratified test set.  Each
row removes or replaces one design component from the full framework.
The geometry-only baseline is excluded from this table because it has no
trainable parameters and does not produce predictions; its spatial
consistency is reported in Table~\ref{tab:spatial_consistency}.  Under
the stratified split, the full framework achieves the lowest MAE,
confirming that the geometric prior contributes predictive signal
rather than merely imposing regularisation.}
\label{tab:ablation}
\begin{tabular}{lcccc}
\toprule
Variant & MAE & RMSE & $R^{2}$ & Pearson $r$ \\
\midrule
Full framework & \textbf{0.417} & \textbf{0.540} & $\mathbf{0.133}$ & 0.567 \\
Without monitoring input & 0.558 & 0.666 & $-0.540$ & $-0.141$ \\
Randomised rock--machine edges & 0.462 & 0.598 & $-0.205$ & 0.438 \\
Without geometric prior ($\beta=0$) & 0.539 & 0.635 & $-0.319$ & $-0.012$ \\
\bottomrule
\end{tabular}
\end{table}

\begin{figure}[htbp]
\centering
\includegraphics[width=0.9\linewidth]{fig4_ablation_study.pdf}
\caption{Structural ablation comparison.  (a)~Prediction metrics (MAE
and $R^{2}$) across four ablation variants; (b)~illustrative spatial
diagnostics generated during figure preparation.  The traceable
checkpoint-based spatial evidence used for interpretation is reported in
Table~\ref{tab:spatial_consistency}.  The geometry-only baseline is
omitted from panel~(a) because it has no trainable parameters and does
not produce predictions.  Under the stratified split, the full framework
achieves the lowest MAE and the only positive $R^{2}$, indicating that
the geometric prior contributes predictive signal rather than merely
imposing regularisation.}
\label{fig:ablation}
\end{figure}

The \emph{No Monitoring} ablation removes the monitoring vector from the GRU
input.  Pearson $r$ collapses from 0.567 to $-0.141$, suggesting that
temporal monitoring dynamics carry essential predictive signal that the
static graph structure alone cannot replace.

The \emph{Randomised Edges} ablation replaces geometry-constrained
rock--machine edges with random edges of the same cardinality.  Pearson $r$
drops from 0.567 to 0.438, suggesting that the spatial structure of
candidate relations---not merely the presence of cross-type
edges---contributes to learning meaningful interaction patterns.

The \emph{No Geometric Prior} ablation sets $\beta=0$.  Under the
stratified split, this variant achieves higher MAE (0.539) and near-zero
Pearson $r$ ($-0.012$) compared with the full framework, indicating that the
geometric prior contributes predictive signal rather than merely
imposing regularisation.  This contrasts with the mileage split, where
removing the prior appeared to improve prediction; the stratified split
isolates the prior's contribution by controlling for geological
heterogeneity between train and test.  Spatial consistency is examined
separately below using post-hoc evidence regenerated from the saved
checkpoint.

The \emph{Geometry-only baseline} replaces all learned edge scores with
the fixed geometric prior $g_{ij}=\exp(-d_{ij}/\tau_{\mathrm{edge}})$
(no trainable parameters).  Because it has no learned forecasting head,
it does not produce predictions and is therefore excluded from
Table~\ref{tab:ablation}; it is, however, retained in the spatial
consistency comparison (Table~\ref{tab:spatial_consistency}) because it
yields a well-defined surface relevance map $C_j$ directly from the
geometric prior.  This baseline serves as a non-learned reference point:
it provides a fixed non-learned spatial reference for evaluating whether
the learned relevance refines or simply mirrors the candidate-edge geometry.

\subsection{Spatial consistency of learned attention}
\label{sec:spatial_consistency}
A central claim of this framework is that geometry-constrained graph
construction produces spatially consistent attention patterns that align
with engineering expectations.  We evaluate this claim using three
quantitative spatial consistency indicators, each targeting a different
aspect of spatial validity.

\paragraph{Attention--geology correlation.}
We quantify whether the mean surface relevance
$\bar{C}(c_k)=|V^m|^{-1}\sum_j C_j(c_k)$ co-varies with the TSP P-wave
velocity interpolated at the same test chainages.  This check is treated
as an external plausibility indicator, not as a required success
criterion: in a response-supervised setting, learned relevance may also
reflect machine operating state and temporal response history.  The
post-hoc checkpoint evidence shows only weak attention--geology alignment
for the learned model (Pearson $r=0.067$, $p=0.837$; Spearman
$r=0.294$, $p=0.354$), whereas the geometry-only prior has a stronger
positive velocity correlation (Pearson $r=0.623$, $p=0.031$).  Thus the
current checkpoint should not be interpreted as proving that low-velocity
geological zones directly receive higher learned relevance.

\paragraph{Spatial autocorrelation of surface relevance.}
Spatially consistent relevance should exhibit positive spatial
autocorrelation: nearby TBM surface nodes should have similar $C_j$ values,
rather than random spatial noise.  We compute Moran's $I$
\citep{moran1950notes} for the surface relevance values on the TBM surface
at each test chainage:
\[
I(c) = \frac{n}{S_0} \cdot
\frac{\sum_i \sum_j w_{ij}\,(C_i(c) - \bar{C}(c))\,(C_j(c) - \bar{C}(c))}
     {\sum_i (C_i(c) - \bar{C}(c))^2},
\]
where $n = |V^m|$ is the number of TBM surface nodes, $w_{ij}$ is the
spatial weight between nodes $i$ and $j$ (binary $k$-nearest-neighbour,
$k=8$, row-standardised), and $S_0 = \sum_i \sum_j w_{ij}$.  Moran's $I$
ranges from approximately $-1$ (perfect dispersion) through $0$ (spatial
randomness) to $+1$ (perfect spatial clustering).  The reported value is
the mean $I(c)$ across all test chainages.  The full framework achieves
Moran's $I=0.925$, indicating strongly clustered surface relevance.  The
geometry-only baseline yields a lower and less stable mean Moran's $I$
(0.170), suggesting that the learned scores are not merely a direct copy
of the fixed geometric prior.

\paragraph{Component-level relevance differentiation.}
If the surface relevance captures meaningful rock--machine interaction
structure, different TBM components (cutterhead, front/middle/tail shield)
should exhibit distinct mean relevance levels, reflecting their different
interaction intensities with the surrounding rock.  We measure the
coefficient of variation (CV) of component-mean relevance across the four
components:
\[
\mathrm{CV} = \frac{\sigma_{\mathrm{comp}}}{\mu_{\mathrm{comp}}},
\quad
\mu_{\mathrm{comp}} = \frac{1}{4}\sum_{k=1}^{4} \bar{C}_k,
\quad
\bar{C}_k = \frac{1}{|V_k^m|}\sum_{j \in V_k^m} C_j,
\]
where $V_k^m$ denotes the set of TBM surface nodes belonging to component
$k$, and $\sigma_{\mathrm{comp}}$ is the standard deviation of the four
component means.  CV is dimensionless and scale-invariant, making it
suitable for comparing differentiation across models with different
absolute relevance magnitudes.  The full framework achieves mean
CV\,$=0.193$, higher than the geometry-only baseline (0.105), indicating
moderate component-level differentiation.

\paragraph{Degree-control check.}
Because rock--machine candidate edges are spatially constrained, nodes
with more incident edges could appear more relevant simply because of
edge count.  We therefore correlate degree-normalised relevance $C_j$
with incident candidate-edge degree $d_j$ across TBM surface nodes.  The
mean Pearson correlation is small ($r=0.043$), suggesting that the
degree-normalised definition of $C_j$ largely removes edge-count inflation.

\begin{table}[htbp]
\centering
\caption{Post-hoc spatial evidence regenerated from the saved checkpoint.
Higher Moran's $I$ and component CV indicate greater spatial structure.
Rel.--Geo.\ $r$ is the Pearson correlation between mean surface relevance
and TSP P-wave velocity across test chainages.  Degree-control $r$
correlates $C_j$ with incident candidate-edge degree and should be close
to zero if relevance is not merely edge-count driven.}
\label{tab:spatial_consistency}
\begin{tabular}{lcccc}
\toprule
Variant & Rel.--Geo.\ $r$ & Moran's $I$ & CV & Degree $r$ \\
\midrule
Full framework & 0.067 & 0.925 & 0.193 & 0.043 \\
Geometry-only baseline & 0.623 & 0.170 & 0.105 & -- \\
\bottomrule
\end{tabular}
\end{table}

These indicators support a narrower but better-evidenced interpretation:
the saved model produces spatially clustered and component-differentiated
surface relevance, and this relevance is not simply explained by the
number of incident candidate edges.  The current evidence does not support
a strong claim that learned relevance is monotonically higher in
low-velocity geological zones; this remains a target for additional
stability and sensitivity experiments.

\subsection{Spatial interpretation outputs}
Figure~\ref{fig:hotspot} presents the response-consistent hotspot maps
projected onto the TBM shield surface at selected chainage positions.
Colour intensity indicates the aggregated edge relevance incident to each
surface node, revealing which shield regions the model identifies as most
strongly coupled to the surrounding rock mass.

\begin{figure}[htbp]
\centering
\includegraphics[width=0.95\linewidth]{fig5_hotspot_maps.pdf}
\caption{Response-consistent surface relevance hotspots on the
unwrapped TBM surface at three representative test chainages.
The colour indicates degree-normalised surface relevance $C_j$,
rescaled within each chainage for visualisation only; different
chainages are therefore not directly comparable in absolute colour
intensity.  Vertical dashed lines separate cutterhead, front shield,
middle shield, and tail shield regions; horizontal dotted lines mark
crown, springline, and invert positions.  (a)~Chainage 41\,m (early
test set); (b)~chainage 45\,m (mid test set); (c)~chainage 48\,m
(late test set).  The maps show where the learned rock--machine
interaction relevance concentrates on the TBM surface.}
\label{fig:hotspot}
\end{figure}

Figure~\ref{fig:chainage_evolution} shows the chainage evolution of
component-wise surface relevance alongside the TSP velocity field and monitoring
responses.  Panel~(a) displays $C_j$ decomposed by TBM component
(cutterhead, front shield, middle shield, tail shield) at each test
chainage, with the cross-component mean overlaid as a dashed line;
panel~(b) shows the TSP P-wave velocity profile with low-velocity
anomalies (below the 25th percentile) highlighted in red; panels~(c)--(g)
show the five monitoring response variables.

\begin{figure}[htbp]
\centering
\includegraphics[width=0.9\linewidth]{fig6_chainage_evolution.pdf}
\caption{Chainage evolution of component-wise surface relevance,
geological context, and monitoring response.  (a)~Surface relevance
$C_j$ decomposed by TBM component (cutterhead, front shield, middle
shield, tail shield; degree-normalised, unstandardised) at each test
chainage, with the cross-component mean overlaid as a dashed line;
(b)~TSP P-wave velocity profile with low-velocity anomalies (below
25th percentile) highlighted in red; (c--g)~monitoring response
variables (advance rate, torque, thrust, penetration, shield pressure).
The co-evolution of component-wise surface relevance with geological
context and monitoring response reveals how interaction emphasis
migrates across TBM components as the TBM advances through different
geological zones.}
\label{fig:chainage_evolution}
\end{figure}

The hotspot maps reveal spatially heterogeneous interaction patterns: the
front shield and cutterhead-proximal region consistently exhibit higher
relevance than the tail shield, consistent with the engineering expectation
that rock--machine interaction is concentrated near the excavation face.
The chainage evolution view shows that interaction relevance shifts across
components as the TBM advances, suggesting that the model captures
geologically meaningful spatial variation in rock--machine coupling.  These
spatial patterns are not recoverable from monitoring-only sequence models,
which lack any spatial representation of the interaction geometry.

Figure~\ref{fig:decision_support} illustrates how the spatial
interpretation outputs \emph{could} inform excavation decision-making by
integrating predicted responses, attention-based signals, and geological
context into a unified spatial view.  This is a conceptual demonstration,
not a validated decision-support tool.

\begin{figure}[htbp]
\centering
\includegraphics[width=0.9\linewidth]{fig7_decision_support.pdf}
\caption{Illustrative decision-support scenario on the test set
(chainage 41--48\,m).  (a)~Predicted vs.\ observed advance rate with
geological scenario annotation (low-velocity zone, transition zone,
normal zone); (b)~surface relevance signal (mean $C_j$) with TSP
velocity overlay on the right axis; the red dashed line marks the
alert threshold (75th percentile of $C_j$), and the shaded red region
highlights chainages where $C_j$ exceeds this threshold;
(c)~composite attention--error indicator combining min--max normalised
surface relevance and normalised prediction error.  This indicator is a
heuristic cue for illustration only; it is not a calibrated risk metric
and has not been validated for operational use.  The spatial reasoning
chain---linking geological context, attention-based interaction
assessment, and prediction error---is \emph{facilitated by} (not
exclusive to) the graph-based spatial representation.  Further
validation on larger datasets is needed before any operational
deployment.}
\label{fig:decision_support}
\end{figure}

\subsection{Robustness and sanity checks}
\label{sec:robustness}
Table~\ref{tab:robustness} summarises the checks that are currently
backed by regenerated evidence files.  The reported values are produced
by reloading the saved checkpoint, rebuilding the same stratified test
set, and recomputing prediction and spatial evidence.  Additional
multi-seed, label-shuffle, and threshold-sensitivity experiments are
defined in the reproducibility protocol but are not reported here until
their output files are generated.

\begin{table}[htbp]
\centering
\caption{Traceable robustness and sanity checks regenerated from the
saved checkpoint.  The evidence source is
\texttt{outputs/evidence/posthoc\_evidence.json}.}
\label{tab:robustness}
\begin{tabular}{llc}
\toprule
Check & Metric & Result \\
\midrule
Checkpoint replay & MAE & 0.417 \\
Checkpoint replay & RMSE & 0.540 \\
Checkpoint replay & $R^2$ & 0.133 \\
Surface clustering & mean Moran's $I$ & 0.925 \\
Component differentiation & mean CV & 0.193 \\
Edge-count control & mean Pearson $r(C_j,d_j)$ & 0.043 \\
Geometry-only reference & mean Moran's $I$ & 0.170 \\
\bottomrule
\end{tabular}
\end{table}

\emph{Checkpoint replay} verifies that the saved graph-sequence model
reproduces the reported prediction metrics when the test set is rebuilt
from the configuration file.  The \emph{edge-count control} rules out the
possibility that higher $C_j$ simply reflects a higher number of incident
candidate edges.  The degree-normalised definition
$C_j = \frac{1}{d_j+\epsilon}\sum_i \mathrm{softplus}(s_{ij}^{rm})$
largely decouples relevance from edge count in the current checkpoint.
Seed stability, label permutation, and candidate-edge threshold
sensitivity remain required follow-up checks before making stronger
claims about robustness across training initialisations or graph
construction parameters.


\section{Discussion}
\subsection{Forward prediction under geological shift}
\label{sec:tradeoff}
The mileage-ordered split evaluates the model on a later continuous tunnel
interval and therefore preserves the geological shift encountered during
excavation.  Under this setting, every model has negative $R^2$.  The proposed
framework obtains MAE 0.435 and is statistically indistinguishable from the
tested baselines.  It should therefore not be presented as a superior
forecasting model on the current dataset.

The ablations reinforce this cautious interpretation.  Removing monitoring
input degrades prediction, and removing the geometric prior produces the worst
learned-model result.  However, randomising the rock--machine edges improves
MAE from 0.435 to 0.413.  Thus the current experiment supports a contribution
from monitoring history and from the prior within the chosen architecture, but
does not demonstrate that the constrained edge topology improves forward
prediction.  The small test set and strong geological shift make this a
particularly demanding evaluation.

\subsection{Spatial evidence and its limits}
The saved-checkpoint diagnostics show positive spatial clustering on the TBM
surface (mean Moran's $I=0.529$).  This establishes that the projected relevance
is spatially structured rather than spatially random.  Component-level
differentiation is weak (mean CV\,$=0.018$), so the current evidence does not
support a strong claim that the model reliably separates cutterhead and shield
interaction roles.

Mean relevance is strongly inversely associated with TSP $V_p$ over the eight
future-test chainages (Pearson $r=-0.988$, $p<0.001$).  Although the direction is
consistent with stronger relevance in lower-velocity ground, this short
contiguous interval also contains strong chainage trends.  The correlation may
therefore reflect common spatial trend rather than a direct geological response.
It is retained as a local diagnostic, not as causal or physical validation.

The degree-control correlation is moderately negative ($r=-0.225$), indicating
that high relevance is not simply produced by more incident candidate edges.
At the same time, degree and relevance are not independent, and additional
topology-shuffle and spatial-trend controls are needed before stronger
interpretability claims are made.

\subsection{Application scenario: spatial reasoning beyond prediction}
\label{sec:application}
The main value currently supported by the framework is representational.  It
connects geological context, monitoring response, candidate rock--machine
relations, and TBM surface locations within one model and can project learned
relevance back to hotspot and chainage views.  Monitoring-only sequence models
cannot produce an equivalent TBM-surface diagnostic because they do not encode
surface nodes or rock--machine relations.

These views remain model diagnostics.  They are not calibrated risk estimates,
contact pressures, jamming probabilities, or prescriptions for machine control.
Operational use would require independent event labels, expert assessment, or
physical measurements that can validate whether the displayed hotspots
correspond to actual interaction states.

\subsection{Boundary and caution}
The case study contains only 44 effective sequence samples and 8 future-test
samples.  The resulting confidence intervals are broad, and the paired tests do
not establish a prediction advantage.  The test interval also differs
geologically from the training interval, so the negative $R^2$ values should be
read as evidence of limited extrapolation rather than ignored as an inconvenient
artifact.

The spatial evidence is checkpoint-based and post-hoc.  Multi-seed stability,
label permutation, candidate-edge threshold sensitivity, spatial detrending,
and cross-project validation remain necessary.  Learned relevance represents
response-consistent model importance under geometric screening, not a direct
measurement of contact force, stress, or engineering risk.

\subsection{Broader GIScience implications}
Despite these limitations, the framework demonstrates a GIScience-relevant
representation strategy for heterogeneous spatial entities under incomplete
observability.  Rock voxels, TBM surface components, candidate relations,
monitoring responses, and projected spatial diagnostics remain linked in one
explicit entity-relation model.

The results also show why spatial graph methods require two distinct evaluation
axes.  Prediction metrics test forward response generalisation, whereas spatial
diagnostics test whether learned relevance has coherent geographic structure.
The current framework provides spatially clustered diagnostics, but neither
predictive superiority nor physical validity is yet established.  Keeping these
claims separate is essential for a defensible interpretation.


\section{Conclusion}
This study builds a geometry-constrained rock--TBM spatial interaction graph
from three inputs: a TSP-derived rock voxel field, an advancing parameterized
TBM surface, and excavation-step-aligned monitoring records. The core method
constructs this graph at each excavation step and aggregates weighted candidate
edges by TBM component into a component
$\times$ excavation-step anomalous interaction field. The resulting field
contains $I_c(t)$ and $\Delta I_c(t)$, can be aligned with monitoring
responses, and retains traceable edge contributions for excavation anomaly
indication, component-indexed readout, and jamming-risk-oriented
interpretation.

The two case studies indicate that the anomalous interaction field organises
rock--machine relations into interpretable component-indexed excavation-step
patterns and, most clearly in the SJLS case, retains residual-consistent
co-variation after simple trend-based null explanations are considered. The
component-level field also provides a step-change readout. In the SJLS case,
$\Delta I_c(t)$ has a large negative association with the step-level
shield-pressure residual change. Relation-support ablations indicate that,
within the present field-readout pairs, the empirical distinction is mainly
between component-indexed candidate-edge aggregation and pooled or global
summaries; the additional numerical effect of distance-normal weighting is
secondary. The evidence supports a spatial field-representation claim rather
than contact-force recovery, calibrated jamming
risk, unique component attribution, or broad predictive dominance.

The proposed field provides a spatial representation basis for future
rock--machine interaction analysis and digital-twin-oriented TBM excavation
analysis. In particular, the component $\times$ excavation-step anomalous
interaction field can be queried alongside monitoring responses to identify
which TBM component is exposed to anomalous rock conditions at which chainage.
It therefore offers a structured input for excavation anomaly indication and
jamming-risk-oriented interpretation in future excavation-analysis workflows.


\section*{Acknowledgements}
The authors thank the project engineers and data managers who supported the
organisation of the TBM monitoring and TSP-derived geological data used in this
study.


\section*{Disclosure statement}
No potential conflict of interest was reported by the author(s).

\section*{Funding}
This work was supported by the National Natural Science Foundation of China
[grant number XXXXXXXX]; and the Fundamental Research Funds for the Central
Universities [grant number XXXXXXXX].

\section*{Data availability statement}
The raw TBM monitoring and TSP data are project data and are not publicly
available because they contain engineering and site information. The materials
that can be shared are the experiment configurations, analysis scripts, figure
scripts, and derived result summaries used to reproduce the reported tables and
figures. These materials are available from the corresponding author upon
reasonable request, subject to project data-use permissions; the raw site data
are not redistributed. To support reproducibility of the spatial graph
construction without disclosing site data, the authors can also provide a
synthetic demonstration package containing a small voxel field, TBM surface
nodes, candidate edge lists, field-calculation scripts, and residual
diagnostic scripts.

\bibliographystyle{plainnat}
\bibliography{references}

\clearpage
\appendix
\section{Exploratory field--response matrix}
\label{app:descriptor_matrix}
Table~\ref{tab:descriptor_residual} reports the full unadjusted
component--response Spearman matrix for all cases, included here for
completeness. These values are exploratory field-readout diagnostics and are
not adjusted for multiple testing.

\begin{table}[htbp]
\centering
\caption{Full unadjusted component--response Spearman matrix for
$I_c(t)$--persistence-residual association. Correlations are exploratory
field-readout diagnostics and are not adjusted for multiple testing.}
\label{tab:descriptor_residual}
\scriptsize
\setlength{\tabcolsep}{3pt}
\begin{tabular}{llrrrrr}
\toprule
Case & Component & Adv. rate & Torque & Thrust & Penetr. & Shield press. \\
\midrule
BSLL h=1 & Cutterhead & 0.238 & -0.286 & 0.048 & -0.095 & 0.095 \\
BSLL h=1 & Front shield & 0.381 & -0.429 & 0.190 & 0.095 & -0.095 \\
BSLL h=1 & Middle shield & -0.048 & 0.262 & 0.452 & -0.167 & 0.143 \\
BSLL h=1 & Tail shield & -0.238 & 0.286 & -0.048 & 0.095 & -0.095 \\
BSLL h=3 & Cutterhead & -0.214 & -0.714 & -0.036 & -0.214 & -0.536 \\
BSLL h=3 & Front shield & 0.286 & -0.179 & 0.143 & 0.250 & -0.750 \\
BSLL h=3 & Middle shield & 0.214 & 0.714 & 0.036 & 0.214 & 0.536 \\
BSLL h=3 & Tail shield & 0.179 & 0.750 & 0.107 & 0.143 & 0.571 \\
SJLS h=3 & Cutterhead & 0.860 & 0.446 & -0.669 & 0.669 & -0.912 \\
SJLS h=3 & Front shield & 0.640 & 0.152 & -0.350 & 0.475 & -0.745 \\
SJLS h=3 & Middle shield & -0.225 & -0.449 & 0.478 & -0.282 & -0.042 \\
SJLS h=3 & Tail shield & -0.846 & -0.324 & 0.733 & -0.811 & 0.784 \\
\bottomrule
\end{tabular}
\end{table}


\end{document}
